%======================================================================
\section{Act VIII --- Give the proof its native types}
\label{sec:haomega}

Everything so far happened inside a deliberately minimal system: one sort,
$\Nat$, and nothing else. That minimality was the point --- it kept the audit
short --- but Act~VI recorded its price as three separate ``documented
compromises,'' and they turn out to be one compromise wearing three coats:
the general Hydra theorem lived outside the fragment because a strategy is a
\emph{function}; Sperner's coloring needed a special symbol because a
coloring is a \emph{function}; and five symbols entered by numeral graph
because their recursions run through an \emph{encoding} of trees into
numbers. All three say: the object language cannot mention $\Nat\to\Nat$.

So we built the same machine again, over \emph{Heyting arithmetic in all
finite types} --- the system modified realizability was invented for
\citep{kreisel1959,troelstra1973}. This act reports what changed, what it
cost, and what the extracted programs look like now. (The repository is
\texttt{modified-realizability-haomega}; the first-order development stays
in-tree as the reference implementation.) One word on the name: throughout
this act, \haomega{} denotes the object theory actually formalized ---
finite-type Heyting arithmetic \emph{extended by} $\mathrm{TI}(\epsz)$ with
its matching recursor, and by the imported case-study primitives below ---
an extension of \haomega{}, not the bare textbook system.

\subsection*{The system, in four displays}

\emph{Types} are the finite types over $\Nat$, with a unit for erased
certificates and two further bases:
\[
  \tau \;::=\; \Nat \;\mid\; \mathbf{1} \;\mid\; \mathsf{ord}
  \;\mid\; \mathsf{hyd} \;\mid\; \tau\to\tau \;\mid\; \tau\times\tau .
\]
$\mathsf{ord}$ is the ordinal notations below $\epsz$ and $\mathsf{hyd}$ is
the hydras --- the two objects the first-order development had to \emph{code}
into $\Nat$. Both are inductive types here, comparison and surgery are pattern
matching, and every certified fact about them is \emph{inherited} from the
coded version by pulling back along the evident encoding, which therefore
survives only inside alignment proofs and appears in no computation at all.
Act~X reports what that buys, and it is measurable.

\emph{Terms} are G\"odel's System~T --- $\lambda$, application, pairing, the
numerals, and a recursor $\mathsf{rec}$ \emph{at every type} --- plus a small
stock of primitives ($+$, the $\epsz$-order $\prec$, and the case-study
symbols $\good$, $\hydra$, $\mathrm{ord}$, \dots, each evaluated by the
first-order development's proven, choice-free value layer --- definable in
System~T in principle, imported as primitives in practice), plus one
genuinely new construct:
\[
  \mathsf{tiRec} : \bigl(\Nat \to (\Nat \to \mathbf{1} \to \tau) \to
  \tau\bigr) \to \Nat \to \tau ,
\]
recursion along $\prec$. Rule and recursor arrive as a \emph{pair}; that
pairing is the subject of Act~IX.

\emph{Formulas} carry their realizer type as an \emph{index},
$\varphi : \mathrm{Formula}\;\Gamma\;\tau$, with the modified-realizability
assignment built into the constructors:
\[
\begin{array}{llcl}
  & s = t,\ \bot &:& \mathbf{1}\\
  & \varphi \land \psi &:& \tau_\varphi \times \tau_\psi\\
  & \varphi \lor \psi &:& \Nat \times (\tau_\varphi \times \tau_\psi)\\
  & \varphi \to \psi &:& \tau_\varphi \to \tau_\psi\\
  & \forall x^{\sigma}.\,\varphi &:& \sigma \to \tau_\varphi\\
  & \exists x^{\sigma}.\,\varphi &:& \sigma \times \tau_\varphi
\end{array}
\]
There is no ambient level, no transport machinery, and --- because
substitution preserves the index \emph{by construction} --- not a single
type-cast anywhere in the extractor. Equality is primitive at every type
(revised from a type-$0$-only first design when Pascal's row-level equations
proved unstatable without it), and \emph{one} Leibniz rule replaces the
fragment's twenty per-symbol congruence schemas.

\emph{Extraction} now lands \emph{in the object language}:
\[
  \extract : \mathrm{Deriv}\;\Delta\;\varphi \;\longrightarrow\;
  \mathrm{Tm}\;(\Delta{+}\Gamma)\;\tau_\varphi .
\]
This is the structural upgrade over Act~III, where the realizer was an
opaque element of the pure-type hierarchy, printable only by walking the
derivation instead. A realizer that is a \emph{term} can be printed,
normalized, emitted as Haskell, and reasoned about --- the repository's
three-view artifact (\texttt{EXTRACTED\_HAOMEGA.md}, rewritten at every
build) renders all thirteen realizers this way, and the displays below are its
three renderings of one certified object each.

\subsection*{Thirteen extracted programs}

\begin{center}\small
\begin{tabular}{ll}
\toprule
theorem & statement (derived in the object theory)\\
\midrule
Goodstein & $\forall m\,\exists t.\ \good(m,t)=0$\\
Kirby--Paris & $\forall h\,\exists t.\ \hydra(h,t)=0$\\
Hercules, $\forall$-strategy & $\forall h\,\forall f^{\Nat\to\Nat}\,\exists t.\ \mathrm{play}(f,h,t)=0$\\
Hercules, any head & $\forall h\,\forall f^{\Nat\to\Nat}\,\forall g^{\Nat\to\Nat}\,\exists t.\ \mathrm{playAt}(g,f,h,t)=0$\\
gcd, full specification & $\forall a\,\forall b\,\exists g.\ g\mid a \land g\mid b \land \forall d.(d\mid a\to d\mid b\to d\mid g)$\\
Pascal mod 2 & $\forall n\,\forall k.\ \mathrm{pas}(n,k)=1 \lor \mathrm{pas}(n,k)=0$\\
Sperner 1D & $\forall n\,\forall c^{\Nat\to\Nat}.\ c\,0=0 \to c\,n=1 \to \exists k.\ k<n \land c\,k\neq c(k{+}1)$\\
Tower of Hanoi & $\forall n\,\exists \mathit{len}\,\exists \mathit{moves}^{\Nat\to\Nat}.\ \dots$\\
Fibonacci & $\forall n\,\exists y.\ y=\fib(n)$\\
Fibonacci at type 2 & $\forall f^{\Nat\to\Nat}\,\exists y.\ y=\fib(f(f\,0))$\\
Goodstein, typed ordinals & $\forall m\,\exists t.\ \good(m,t)=0$ \ (again, on $\mathsf{ord}$)\\
Kirby--Paris, typed trees & $\forall h^{\mathsf{hyd}}\,\exists t.\ \mathrm{dead}(\mathrm{play}(h,t))=0$\\
Hercules any head, typed trees & $\forall h^{\mathsf{hyd}}\forall f\forall g\,\exists t.\ \mathrm{dead}(\mathrm{playAt}(g,f,h,t))=0$\\
\bottomrule
\end{tabular}
\end{center}

Four of these the fragment could not even \emph{state}: the two Hercules
theorems and Sperner quantify over functions, and the type-2 Fibonacci is
where the continuity theorem finally has content ($\mathsf{Continuous2}$ of
the extracted functional, with an explicit computable modulus
$\max(f\,0, f(f\,0))+1$). The any-head Hercules closes what the audit had
recorded as the remaining gap: a computable surgery move ($\mathsf{hcutAt}$,
the leftmost-head move with the head position as an argument) whose descent
is the first-order \texttt{play\_descends} read on codes, so quantifying
the head strategy costs one primitive and one schema. One scope statement
survives, verbatim from the audit: independence from \textsc{pa} is
formalized nowhere.

\subsection*{Reading an extracted program, three ways}

Fibonacci first, because it fits on a line. The raw extracted object ---
the realizer itself, with its erased certificate visible as $\star$:
\begin{quote}\ttfamily\small
($\lambda$x0.\ (($\lambda$x1.\ fst (($\lambda$x2.\ rec[$\langle$0, S 0$\rangle$
$\mid$ ($\lambda$x3.\ ($\lambda$x4.\ (($\lambda$x5.\ $\langle$snd x5, (fst x5 + snd
x5)$\rangle$) x4))) $\mid$ x2]) x1)) x0), $\star$)
\end{quote}
Collapsed --- contentless parts elided, what remains is exactly the
computational content --- it is the same term minus the trailing
$\star$-component: iterate $\langle a,b\rangle \mapsto \langle b,
a{+}b\rangle$ from $\langle 0,1\rangle$, take the first projection. And as
generated Haskell (see the certification note below):
\begin{quote}\ttfamily\small
(\char`\\ x0 -> (((\char`\\ x1 -> (fst ((\char`\\ x2 -> (natRec (0, 1)\\
\hspace*{1.2em}(\char`\\ x3 -> (\char`\\ x4 -> ((\char`\\ x5 -> ((snd x5), ((fst x5) + (snd x5)))) x4)))\\
\hspace*{1.2em}x2)) x1))) x0), ())
\end{quote}

Now the one worth the whole apparatus: Goodstein's collapsed program,
abridged to its skeleton with the bureaucracy elided ($\ldots$), because the
mathematics is \emph{visible in it}:
\begin{quote}\ttfamily\small
($\lambda$m.\ (\dots tiRec[\,($\lambda$x.\ ($\lambda$ih.\ (\dots\\
\hspace*{1.2em}rec[\,$\langle$s,\ \dots$\rangle$
  \hfill\normalfont\itshape --- eqDec hits $0$: return the stop time $s$\\
\hspace*{1.2em}$\mid$ \dots (ih\ ord(S S S s,\ good(m, S s))) \dots]
  \hfill\normalfont\itshape --- else recurse at the descended ordinal\\
\dots))\,] \ ord(S S 0,\ m)\,)
  \hfill\normalfont\itshape --- start at $\mathrm{ord}(2,m)$
\end{quote}
The transfinite recursor, the zero test on $\good(m,s)$, and the $\epsz$
descent $\mathrm{ord}(s{+}3, \good(m,s{+}1)) \prec \mathrm{ord}(s{+}2,
\good(m,s))$ --- the entire proof structure, readable in the program it
became. The extracted stopping-time function returns the published values
$[0,1,3,5]$ for $m=0..3$, checked at every build together with the
certificate $\good(m, \mathrm{stop}(m))=0$; the first-order extract, walled
behind its ambient tower, could not evaluate past $m=1$.

\section{Act IX --- Matched principles: a rule and its recursor}

Adding $\mathrm{TI}(\epsz)$ to the object theory is not a free move, and it is
worth being exact about why, because it is the one place where this
development strengthens the logic rather than merely retyping it.

System~T defines exactly the provably-total functions of \textsc{pa}. Goodstein's
stopping time is not among them. So a proof rule strong enough to derive
Goodstein's theorem cannot have its computational content expressed by the
term language as it stands: extraction would reach the rule and have nothing
to emit. There are two ways out, and only one of them is honest. The
dishonest one is to let the \emph{evaluator} do the work --- to define the
extracted program by an unbounded search, or by a recursion whose termination
the metatheory knows but the object language does not. That buys running
programs at the price of the thing being demonstrated, since the program is
then no longer read off the proof.

The route taken here is the other one. The rule arrives with a term former:
\[
  \mathrm{TI}(\epsz)
  \quad\longleftrightarrow\quad
  \mathsf{tiRec} : \bigl(\Nat \to (\Nat \to \mathbf{1} \to \tau) \to
  \tau\bigr) \to \Nat \to \tau .
\]
Read the type: at $x$, given the values at every $y$ \emph{together with a
realizer of $y \prec x$}, produce the value at $x$. That is the rule's premise
transcribed. And because the order premise is an equation, its realizer is
contentless --- the $\mathbf{1}$ argument --- so the recursor cannot receive
evidence that a recursive call is legal, and must re-decide $\prec$ itself and
fall back when the test fails. The proof rule and the program construct are
the same shape, and neither is doing work the other does not license.

We can state the principle generally, because the development applies it
twice. When the object theory is strengthened by a proof principle whose
content the term language cannot already express, the strengthening must come
paired with a term former whose type is the principle's premise read as a
type. Both instances here obey it: $\mathrm{TI}(\epsz)$ with $\mathsf{tiRec}$
over coded ordinals, and its twin over the structural notations of
$\mathsf{ord}$.

What does \emph{not} obey it, and should be named plainly, is the second kind
of extension: the case-study symbols. $\good$, $\hydra$, $\mathrm{ord}$ and
the rest enter as primitives evaluated by the first-order development's
verified value layers. They are definable in System~T in principle; they are
imported in practice, because defining hereditary base representation and tree
surgery inside the object language is a project of its own and not this one.
That is a convenience, not a principle, and it is why the name \haomega{} is
qualified every time it appears here: the object theory formalized is
finite-type Heyting arithmetic extended by a matched pair \emph{and} by a
stock of verified primitives. A reader who objects that this is not the bare
textbook system is right, and the more precise description is the more
interesting one --- a framework for finite-type realizability that is
\emph{extensible by matched proof/program principles}, with one worked
extension of genuine proof-theoretic strength.

%======================================================================
\section{Act X --- Three experiments}

Acts~I--VIII establish that the pipeline works, twice, over two object
languages. What makes the second system worth building is that having two of
them turns features of the setup into \emph{variables}. Three can be moved one
at a time, with the extraction mechanism held fixed, and in each case the
extracted program registers the change.

\subsection*{Varying the proof: the program is the proof's fingerprint}

Sperner 1D was proved here by the same forward-scan invariant as in Act~IV:
$P(m) := c\,m = 0 \lor \exists k{<}m.\ c\,k \neq c(k{+}1)$. Same theorem,
same invariant --- yet the extracted searches \emph{differ measurably}. The
first-order proof consults the invariant first and decides $c(m{+}1)$
second; its program returns the \emph{first} crossing. The \haomega{} proof
decides $c(m{+}1) = 0$ first, so a fresh zero re-enters the left disjunct
and discards any crossing already found; its program returns the
\emph{last}: on $[0,1,0,1]$ the first-order extract answers $0$, this one
answers $2$ --- both certified crossings, at every input, by soundness.
Nothing in Acts~I--VII showed this cleanly: two proofs of one theorem,
distinguishable by running their extracts. The extraction map is faithful
to proof structure --- which is the thesis of Act~VII, here caught in the
act.

\subsection*{Varying the representation}

The first experiment varied the proof. The second holds the proof fixed ---
line for line --- and varies how the objects it reasons about are
\emph{represented}. \texttt{GoodsteinTyped.lean} proves
$\forall m\,\exists t.\ \good(m,t)=0$ a second time, by transfinite induction
on structural notations instead of codes --- the same derivation line for
line, with three symbols swapped ($\mathrm{ord}\mapsto\mathrm{ord^o}$,
$\prec\mapsto\prec^o$, $\mathsf{tiRec}\mapsto\mathsf{tiRec^o}$). The
extracted program contains \emph{no coded ordinal}: zero occurrences of the
coded $\mathrm{ord}$, $\prec$ and $\mathsf{tiRec}$, and the two programs are
guarded equal at every input the build evaluates.

What that buys is measured, and it is a claim about \emph{size}, not speed:

\begin{center}\small
\begin{tabular}{rrr}
\toprule
$n$ & code (decimal digits) & tree (nodes)\\
\midrule
$10$ & $3$ & $13$\\
$1{,}000$ & $1{,}412$ & $67$\\
$100{,}000$ & $1{,}782$ & $83$\\
\bottomrule
\end{tabular}
\end{center}

No reproducible running-time difference between the two extracted programs
was found, and none is claimed --- Lean's naturals are \textsc{gmp}-backed,
so arithmetic on a $1{,}782$-digit code is cheap. What the typed layer
removes is the encoding itself, and with it the decode and normal-form side
conditions that the first-order ordinal proofs carry throughout.

For the \emph{hydras}, though, the encoding was not merely tidy-up: it was
the measured bottleneck. \texttt{HydraTyped.lean} gives trees the base type
$\mathsf{hyd}$, with the Kirby--Paris ordinal landing in $\mathsf{ord}$, so
a battle state and its measure are both structural; the descent is the
first-order \texttt{play\_descends} applied to trees with no encode/decode
round trip, and the whole layer adds \emph{one} rule where the coded version
needs three. \texttt{HydraTree.lean} then re-proves the theorem, and its
extract contains zero coded operations. The consequence is the sharpest
measurement in this act: at the hydra whose published Kirby--Paris battle
length is $37$, the coded extract overflows the interpreter without taking a
single step, and the tree extract returns $37$. The wall was never the
mathematics, the proof, or the type discipline --- it was the pairing
function.

With the general Hercules game re-proved on trees as well
(\texttt{HerculesTree.lean}), the de-coding programme is complete: every
case study that once coded an object into $\Nat$ has a typed twin, and the
coded modules survive only as the baselines their twins are measured
against. That last port is also the cheapest, and instructively so --- at
the value level it cost \emph{nothing}, because the any-head move and its
descent were already functions and theorems about trees. What had been
coded was never the mathematics; it was only the interface the one-sorted
object language forced on it.

\subsection*{Varying the logical strength}

The third axis is the one Part~I already exercised, but it is only visible as
a \emph{variable} now that the other two have been moved. Take the same
theorem shape --- $\forall m\,\exists t.\ \dots$ --- and vary how much
induction the proof is allowed. An ordinary induction extracts to
$\mathsf{rec}$, recursion along the successor. A proof that needs
$\mathrm{TI}(\epsz)$ extracts to $\mathsf{tiRec}$, recursion along $\prec$.
The printed Goodstein skeleton above shows it directly: the outermost
recursion is not over $\Nat$ but over the ordinal, and the induction
hypothesis appears as the recursive call at the descended ordinal. Raising the
proof-theoretic strength does not merely license more theorems; it changes the
shape of the recursion the extractor emits, and the change is legible in the
program text.

\subsection*{The comparison, tabulated}

Setting the two systems side by side --- same extraction mechanism, same
theorems, different object language --- gives the following. Every row is
either a statement about what can be \emph{said}, or a measurement.

\begin{center}\small
\begin{tabular}{@{}lp{0.29\linewidth}p{0.29\linewidth}@{}}
\toprule
 & first-order (coded) & finite type (structural)\\
\midrule
quantify a strategy & not statable & a variable $f^{\Nat\to\Nat}$\\
a colouring & the \code{look} primitive & a function $\Nat\to\Nat$\\
ordinal notations & naturals, via a pairing & base type $\mathsf{ord}$\\
hydra trees & naturals, via a pairing & base type $\mathsf{hyd}$\\
binders & two capture-avoidance devices & intrinsically typed; capture impossible\\
congruences & ${\sim}20$ per-symbol schemas & one Leibniz rule\\
the realizer & element of the pure-type tower & a term of the object language\\
ambient levels & up to $41$ (gcd) & none\\
inference rules & $76$ & $45$\\
\midrule
gcd extract & certified; evaluates at no input & evaluates\\
Tower of Hanoi & reaches $n = 4$ & reaches $n = 10$\\
Kirby--Paris, the $37$-battle & overflows the evaluator & returns $37$\\
\bottomrule
\end{tabular}
\end{center}

The rows above the rule are about expressiveness and bookkeeping; the three
below it are the ones that cost programs. And the honest reading of the
division is that the first block \emph{causes} the second: a strategy that
cannot be quantified forces a metatheorem, an ordinal that must be coded
forces arithmetic on thousand-digit integers, and the levels that exist only
to move realizers between tiers are what stop \texttt{gcdWitness} from
evaluating at any input at all.

%======================================================================
\section{Act XI --- What the extracted programs satisfy}

Act~III proved that every extracted program is continuous, and the proof was
real, but in the fragment it had nothing to bite on. Continuity constrains how
a program may inspect an \emph{infinite} input; every theorem the fragment can
state quantifies over numbers, so every extract takes a number, and the
constraint is satisfied vacuously. The machinery was pre-paid infrastructure.

With function variables it becomes a claim with content. The type-2 Fibonacci
theorem $\forall f^{\Nat\to\Nat}\exists y.\ y = \fib(f(f\,0))$ extracts to a
functional
\[
  \mathrm{hiProgram} : (\Nat\to\Nat) \to \Nat ,
\]
whose input is an infinite object. The continuity theorem now says something
one can check: the program consults only finitely much of $f$. For this
program the finite part is explicit and computable ---
\[
  \mathrm{hiModulus}(f) \;=\; \max\bigl(f\,0,\ f(f\,0)\bigr) + 1 ,
\]
with the theorem that any $g$ agreeing with $f$ below that bound gives the
same answer. It is also exhibited as an element of $\CtQ 2$, with a type-$1$
associate: the type-2 functional is represented by an ordinary function on
numbers.

The sharper evidence that the theorem is not vacuous is negative. Consider
\[
  \mathrm{notAllZero}(\alpha) \;=\;
    \begin{cases} 0 & \text{if } \alpha\,n = 0 \text{ for every } n\\
                  1 & \text{otherwise,}\end{cases}
\]
a perfectly well-defined functional of the same type. Deciding it requires
reading all of $\alpha$, so it is not continuous --- and the development
proves the consequence: \emph{no derivation extracts to it}. There is an
inhabitant of the type on the far side of a line the logic cannot cross, and
the line is drawn by a theorem rather than by inspection.

How the proof goes is worth one paragraph, because it is the same device that
reappears in Act~XII. The vendored Kleene--Kreisel development is indexed by
pure-type \emph{level} and has neither products nor general arrows, while
realizers here live at arbitrary finite types. Rather than generalize that
library, the proof carries an oracle-parameterized logical relation
$\mathrm{Tracked}\,\tau\,X$: continuity at base type is the vendored notion,
and every other type gets a clause. No associates are ever constructed, so the
pure-type tower's shape never has to be matched --- and the induction runs
over the constructors of the term language, one case each, where the
first-order version needed a preservation lemma for each of some forty
extraction combinators.

%======================================================================
\section{Act XII --- From derivation to running code}

The Haskell renderings carried a disclaimer for most of this development:
\emph{uncertified translation; the certified artifact is the term}. That
disclaimer is now a theorem plus a short, explicit trusted base.

Certifying emitted code against \textsc{ghc} is not available to us --- it
would need a formal semantics of Haskell and a proof that \textsc{ghc}
implements it. What is available is what a verified compiler actually proves:
correctness against a formal semantics of the \emph{target}. So the target
fragment gets a syntax $\mathrm{HsTm}$ (untyped, since emission erases every
distinction $\mathrm{Ty}$ makes), a big-step evaluation relation with closures, and a
translation $\mathrm{hsOf}$. Agreement across a type erasure cannot be an
equation, so it is a logical relation $\mathrm{Rel}\;\tau\;x\;v$ --- an
equation at base types, closure under application at arrows. That is the same
device the continuity theorem uses, put to a second purpose:
\[
  \mathrm{hsOf\_correct} :\ \text{related environments} \Rightarrow
  \exists v.\ \mathrm{HsEval}\ (\mathrm{hsOf}\ t)\ \rho\ v \ \wedge\
  \mathrm{Rel}\ \tau\ (t.\mathrm{eval}\ e)\ v .
\]
The emitter now \emph{is} $\mathrm{hsPrint} \circ \mathrm{hsOf}$, so the
strings in the artifact come from the syntax tree the theorem is about; the
regenerated artifact is byte-identical, which is the evidence that routing it
through the certified translation changed nothing.

What remains trusted is worth stating precisely, because it is short: the
hand-written prelude, whose meaning the model \emph{defines} to be the
corresponding Lean primitives; the printer, and \textsc{ghc}'s parse of its
output; and the seven programs that use $\mathrm{TI}(\epsz)$, for which the
emitter has never produced running code --- $\mathrm{hsOf}$ sends
$\mathsf{tiRec}$ to a constructor with no evaluation rule, faithfully
modelling the \texttt{error} the prelude emits. Six of the thirteen programs
are covered, and which six is checked at every build.

%======================================================================
\section{Act XIII --- The ledger, and what the two systems show}

\subsection*{The audit, rerun}

The footprints, from the machine, for the \haomega{} library:
\begin{center}\small
\begin{tabular}{ll}
\toprule
$\extract$, the Fibonacci realizer & \emph{no axioms}\\
continuity; every derivation; every running extract & \texttt{[propext, Quot.sound]}\\
$\mathsf{soundness}$ & \texttt{[propext, Classical.choice, Quot.sound]}\\
\bottomrule
\end{tabular}
\end{center}
Choice enters $\mathsf{soundness}$ exactly where it did in Act~V: through the
value-layer \emph{theorem proofs} discharging the ordinal-descent schemas,
while the value-layer definitions --- everything that runs --- stay
choice-free. The verdict of Part~I therefore survives the upgrade unchanged:
every extracted program that runs depends on two inert axioms and nothing
classical, and the kernel will reprint that on demand.

\subsection*{What is not claimed}

The ledger is more useful for what it refuses than for what it asserts, so
here it is in one place. Independence from \textsc{pa} is formalized nowhere:
Goodstein and Kirby--Paris are proved here as termination theorems, and that
\textsc{pa} cannot prove them is not claimed. The typed ordinal layer is a
result about representation \emph{size}; no reproducible running-time
improvement was found for the two Goodstein extracts, and none is asserted.
The certified translation covers six of the thirteen programs --- the seven
that use $\mathrm{TI}(\epsz)$ are not covered, because the target has no
transfinite recursor and the emitter has never produced running code for
them. The case-study symbols are imported primitives, not System~T
definitions. And the coded modules retained as baselines are still coded:
their hydra extracts still overflow the evaluator at codes $4$ and $7$, which
is the measurement their typed twins are measured against, not a defect that
was repaired in place.

None of these is a footnote added under pressure. Each is a claim the
development at some point stated more strongly and then walked back when the
machine, or a measurement, declined to support it.

\subsection*{The thesis}

Part~I answered the question the Prologue asked: a proof that needs more than
Peano arithmetic becomes a program that needs less than a whiff of choice, and
we know it does because the machine was asked. That is a result about
\emph{extraction} --- the mechanism works, at strength, honestly.

Part~II answers a question Part~I could not pose. Holding the extraction
mechanism fixed and moving one feature of the setup at a time, the extracted
program registers each move. Change the proof and the program changes
strategy: two derivations of Sperner's lemma, same theorem and same invariant,
return different crossings. Change the representation and the program changes
what it can finish: the same derivation, line for line, computes a $37$-step
hydra battle from trees that it cannot begin from codes. Change the logical
strength and the recursion changes shape: $\mathsf{rec}$ becomes
$\mathsf{tiRec}$, and the ordinal descent is visible in the program text.
Change the type and semantic constraints appear that were vacuous before:
continuity acquires content exactly when a realizer first takes a function as
input, and the functionals it excludes can be named.

Put together, the two parts say something narrower and more useful than
``proofs have computational content.'' They say that the content is
\emph{sensitive} --- to how the theorem was proved, to how its objects were
represented, and to how much induction it was allowed --- and that with the
extraction mechanism held fixed and the axiom ledger machine-checked, those
sensitivities can be measured rather than argued about. What had been coded
was never the mathematics; it was only the interface the one-sorted object
language forced on it, and the cost of that interface is now a number rather
than an intuition.

That is how a proof becomes a program: not by a grand mechanism, but by a
small one applied with discipline, checked by something that cannot be talked
past --- and, once there are two of them to compare, by a machine in which the
consequences of proving things one way rather than another can be run.
